% Generated draft: requires mathematical and source review.
\documentclass[11pt]{article}
\usepackage[margin=25mm]{geometry}
\usepackage{fontspec}
\setmainfont{Latin Modern Roman}
\usepackage{amsmath,amssymb,mathtools,unicode-math}
\setmathfont{Latin Modern Math}
\usepackage{enumitem,needspace,xurl,hyperref}
\setlength{\parindent}{0pt}
\setlength{\parskip}{5pt}
\setlength{\emergencystretch}{3em}
\newcommand{\N}{\mathbb{N}}
\newcommand{\Z}{\mathbb{Z}}
\newcommand{\Q}{\mathbb{Q}}
\newcommand{\R}{\mathbb{R}}
\newcommand{\C}{\mathbb{C}}
\newcommand{\F}{\mathbb{F}}
\newcommand{\A}{\mathbb{A}}
\newcommand{\PP}{\mathbb{P}}
\newcommand{\E}{\mathbb{E}}
\newcommand{\eps}{\varepsilon}
\newcommand{\dd}{\,\mathrm{d}}
\begin{document}
\section*{UnsolvedMath ID: 30002064}
\textbf{Source problem code:} catalogue entry
\subsection{Definitions and mathematical statement}
The supplied record names the following the Matveev--Jaco--Rubinstein lens-space complexity conjecture. To make its statement precise, this draft uses Matveev complexity and regular continued fractions, as defined below. These conventions are customary interpretations of the terminology; the supplied record does not explicitly define them, so they require confirmation against the cited source.

Let \(p,q\in\Z\) satisfy
\[
 p>q>0,\qquad p>3,\qquad \gcd(p,q)=1.
\]
Using a standard lens-space convention, let
\[
 S^3=\{(z_1,z_2)\in\C^2:|z_1|^2+|z_2|^2=1\},
 \qquad \zeta=\exp(2\pi i/p),
\]
and define \(L(p,q)\) to be the quotient of \(S^3\) by the cyclic action generated by
\[
 (z_1,z_2)\longmapsto(\zeta z_1,\zeta^q z_2).
\]
The source does not specify its lens-space convention or whether orientation is retained; the complexity used here concerns the underlying PL manifold.

For a closed connected PL \(3\)-manifold \(M\), a spine is a compact polyhedron \(P\) of dimension at most \(2\) onto which \(M\setminus\operatorname{int}(B^3)\) collapses, where \(B^3\subset M\) is an embedded PL ball. Such a spine is almost simple if the link of every point of \(P\) embeds in the complete graph \(K_4\), equivalently the \(1\)-skeleton of a tetrahedron. A true vertex is a point whose link is homeomorphic to \(K_4\). Define the Matveev complexity \(c(M)\) as the minimum number of true vertices among all almost simple spines of \(M\).

Write the finite regular continued fraction of \(p/q\) as
\[
 \frac pq=[a_0;a_1,\ldots,a_n]
 =a_0+\cfrac{1}{a_1+\cfrac{1}{\ddots+\cfrac{1}{a_n}}},
\]
where \(n\geq0\), \(a_0\geq1\), \(a_i\geq1\) for \(1\leq i\leq n\), and \(a_n\geq2\) when \(n\geq1\). When \(n=0\), the expression means \([a_0]=a_0\). Set
\[
 E(p,q)=\sum_{i=0}^{n}a_i.
\]
Here the source's phrase ``partial denominators'' is interpreted as all the coefficients in this regular continued fraction, including \(a_0\). The source does not explicitly distinguish regular continued fractions from other continued-fraction conventions.

The conjectured statement, under these conventions, is
\[
 \text{for every }p,q\in\Z\text{ with }p>q>0,\ p>3,\ \gcd(p,q)=1,
 \qquad c\bigl(L(p,q)\bigr)=E(p,q)-3.
\]
The supplied record labels the problem partially solved and reports results for infinite families and many small cases. Those are supplied literature-status claims, not independently verified conclusions of this draft.
\subsection{Short English statement}
For every coprime pair of integers \(p>q>0\) with \(p>3\), the conjecture asserts that the Matveev complexity of \(L(p,q)\) equals the sum of the coefficients of the regular continued fraction of \(p/q\), minus three. The precise complexity, lens-space, and continued-fraction conventions require source confirmation. The supplied status is partially solved.
\subsection{Sources}
\begin{itemize}
\item Local UnsolvedMath registry: permanent ID \(30002064\), title ``Matveev--Jaco--Rubinstein Lens-Space Complexity Conjecture,'' source problem code \texttt{catalogue entry}, in the Oberwolfach Reports Open Problems collection. The record gives proposed year \(2012\), names no individual proposer, and attributes the conjecture by its Matveev--Jaco--Rubinstein name.
\item Supplied primary report citation and literature-source label ``original OWR report'': \emph{Triangulations (2013)}, Oberwolfach Reports, pp.\ 1405--1486, DOI \texttt{10.4171/owr/2012/24}. \url{https://doi.org/10.4171/owr/2012/24}. The supplied background describes these reports as records of mathematical workshops and their open-problem discussions. The citation's year \(2013\), the proposed year \(2012\), and the DOI containing \(2012\) are retained as supplied; their relationship has not been checked.
\item Supplied literature source, labelled ``minimal lens triangulations'': \url{https://arxiv.org/abs/0805.2425}.
\item Supplied literature source, labelled ``stable-complexity literature'': \url{https://doi.org/10.1112/jtopol/jts026}.
\item Supplied literature assessment dated exactly \texttt{2026-08-21}: the formula is reported proved for infinite families of lens spaces and many small cases, while the general exact-complexity conjecture for all lens spaces is reported unresolved. This is supplied dated curation derived from the OWR material, not an independently verified literature assessment. The source itself requests independent verification of time-sensitive status and cited claims.
\item The registry labels the statement \texttt{corrected\_verified} and reports a manual semantic, typographic, and LaTeX review of its cleanup. This is supplied editorial metadata, not a mathematical verification performed for this draft.
\end{itemize}
\end{document}
